\SourceLeaf{638}
\begin{center}ZUM GEBRAUCH DES KOMPARATORS.\end{center}
\begin{center}[1.]\end{center}
Drei Mikroskope, die in annähernd gleichen Abständen aufgestellt sind, werden nacheinander über die Abschnitte eines Maßstabs gebracht, die diesen Abständen annähernd entsprechen.

Bei diesen aufeinanderfolgenden Versuchen gehen die Teilstriche des Maßstabs um die folgenden Beträge über die Sehlinien der Mikroskope hinaus
\[
\begin{array}{ccc}
a^0&b^0&c^0\\
a'&b'&c'\\
a''&b''&c''\\
a'''&b'''&c'''\\
&\text{etc.}&
\end{array}
\]
Mit $x^0,y^0,z^0;\ x',y',z';\ x''$ usw. bezeichnen wir die Fehler dieser Größen. Setzt man
\[
\begin{aligned}
a'-a''-b^0+b''+c^0-c'&=e^0\\
a''-a'''-b'+b'''+c'-c''&=e'
\end{aligned}
\]
usw., so lauten die Fehlergleichungen
\[
\begin{aligned}
x'-x''-y^0+y''+z^0-z'&=e^0\\
x''-x'''-y'+y'''+z'-z''&=e'\ \text{usw.}
\end{aligned}
\]
Danach haben die plausibelsten Werte der Fehler die Form
\[
\begin{array}{lll}
x^0=0&y^0=-h^0&z^0=h^0\\
x'=h^0&y'=-h'&z'=h'-h^0\\
x''=h'-h^0&y''=-h''+h^0&z''=h''-h'\\
x'''=h''-h'&y'''=-h'''+h'&z'''=h'''-h''\\
\text{usw.}&&
\end{array}
\]
\SourceLeaf{639}
und die Hilfsgrößen $h^0,h'$ usw. hängen von den Gleichungen ab
\[
\begin{aligned}
6h^0-2h'-h''&=e^0\\
-2h^0+6h'-2h''-h'''&=e'\\
-h^0-2h'+6h''-2h'''-h''''&=e''\\
-h'-2h''+6h'''-2h''''-h^{\mathrm V}&=e'''\\
-h''-2h'''+6h''''-2h^{\mathrm V}-h^{\mathrm {VI}}&=e''''\\
\text{etc.}&
\end{aligned}
\]

\par\Needspace{5\baselineskip}
\begin{center}[2.]\end{center}
Anordnung der Längenvergleiche, um einen Abschnitt eines geteilten Maßstabs in Teilen des gesamten Maßstabs zu bestimmen.

\textit{Beispiel.} Die Teilstriche des Maßstabs waren mit Zahlen von 0 bis 840 bezeichnet. Der Abschnitt von 0 bis 87 sollte in Teilen des gesamten Maßstabs von 0 bis 840 bestimmt werden. — Als Abkürzung soll
\[
0\ .\ 87\ .\ 174=\alpha
\]
bedeuten, dass sich beim Vergleich der Länge von 0 bis 87 mit der Länge von 87 bis 174 ergeben hat, dass die erste um $\alpha$ Mikrometerteilungen größer ist als die zweite, und so fort.
\begin{center}Anordnung:\end{center}
\[
\begin{array}{rcrcrcl}
0&.&87&.&174&=&\alpha\\
0&.&174&.&348&=&\beta\\
0&.&348&.&696&=&\gamma\\
840&.&696&.&552&=&\delta\\
840&.&552&.&264&=&\varepsilon\\
0&.&264&.&528&=&\zeta\\
840&.&528&.&216&=&\eta\\
0&.&216&.&432&=&\theta\\
840&.&432&.&24&=&\iota\\
0&.&24&.&48&=&\kappa\\
0&.&48&.&96&=&\lambda\\
0&.&96&.&192&=&\mu\\
0&.&192&.&384&=&\nu\\
0&.&384&.&768&=&\pi\\
840&.&768&.&696&=&\rho
\end{array}
\]

\SourceLeaf{640}
\textit{Berechnung.} Bezeichnet man die gesamte Länge des Maßstabs mit $840'$, die gesuchte Länge (von 0 bis 87) mit $87'$, die Länge von 87 bis 174 mit $174'-87'$ usw., so erhält man die folgenden Gleichungen:
\[
\begin{aligned}
2\ .\ 87'-\alpha&=174'\\
2\ .\ 174'-\beta&=348'\\
2\ .\ 348'-\gamma&=696'\\
2\ .\ 696'-840'+\delta&=552'\\
2\ .\ 552'-840'+\varepsilon&=264'\\
2\ .\ 264'-\zeta&=528'\\
2\ .\ 528'-840'+\eta&=216'\\
2\ .\ 216'-\theta&=432'\\
2\ .\ 432'-840'+\iota&=24'\\
2\ .\ 24'-\kappa&=48'\\
2\ .\ 48'-\lambda&=96'\\
2\ .\ 96'-\mu&=192'\\
2\ .\ 192'-\nu&=384'\\
2\ .\ 384'-\pi&=768'\\
2\ .\ 768'-840'+\rho&=696'
\end{aligned}
\]
Daraus ergibt sich
\[
\begin{aligned}
696'&=-840'+\rho-64.840'+64\iota-256.840'+256\eta-1024.840'\\
&\quad-2\pi-128\theta-512\zeta+1024\varepsilon\\
&\quad-4\nu-2048.840'\\
&\quad-8\mu+2048\delta\\
&\quad-16\lambda+4096.696'\\
&\quad-32\kappa
\end{aligned}
\]
Somit gilt, wenn $840'=L$, $696'=y$, $87'=z$,
\SourceLeaf{641}
\[
\begin{aligned}
3393L&=4095y+\rho\\
&\qquad-2\pi\\
&\qquad-4\nu\\
&\qquad-8\mu\\
&\qquad-16\lambda\\
&\qquad-32\kappa\\
&\qquad+64\iota\\
&\qquad-128\theta\\
&\qquad+256\eta\\
&\qquad-512\zeta\\
&\qquad+1024\varepsilon\\
&\qquad+2048\delta
\end{aligned}
\]
\[
z=\tfrac18y+\tfrac12\alpha+\tfrac14\beta+\tfrac18\gamma
\]
Danach erhält man den gesuchten Wert von $z=87'$ in Teilen des gesamten Maßstabs $L=840'$, wenn das Mikrometer geprüft ist und dadurch feststeht, dass der Wert seiner Teilungen gleich ist und in Teilen des Abstands zweier nach der obigen Anordnung eingestellter Teilstriche des Maßstabs (zum Beispiel in Teilen der Länge von 87 bis 96) angegeben ist.

